\documentclass[11pt,a4paper]{article}

\usepackage[utf8]{inputenc}
\IfFileExists{lmodern.sty}{\usepackage[T1]{fontenc}\usepackage{lmodern}}{}
\usepackage[margin=2.5cm]{geometry}
\usepackage{amsmath,amssymb,amsthm}
\usepackage{bm}
\usepackage{booktabs}
\usepackage{graphicx}
\usepackage{float}
\usepackage{caption}
\usepackage{xcolor}
\usepackage[colorlinks=true,linkcolor=blue!60!black,citecolor=blue!60!black,urlcolor=blue!60!black]{hyperref}
\usepackage{enumitem}
\setlist{itemsep=2pt,topsep=4pt}

\newcommand{\E}{\mathbb{E}}
\newcommand{\Var}{\operatorname{Var}}
\newcommand{\CVaR}{\operatorname{CVaR}}
\newcommand{\VaR}{\operatorname{VaR}}
\newcommand{\PnL}{\mathrm{PnL}}

% Include a figure only if the file exists (keeps the document compilable without figures)
\newcommand{\maybefigure}[3]{%
  \IfFileExists{#1}{%
    \begin{figure}[H]
      \centering
      \includegraphics[width=#2\textwidth]{#1}
      \caption{#3}
    \end{figure}%
  }{%
    \begin{center}\fbox{\parbox{0.8\textwidth}{\centering\small Figure \texttt{\detokenize{#1}} not found.\\ #3}}\end{center}%
  }%
}

\title{\textbf{Deep Hedging under Transaction Costs}\\[4pt]
\large CVaR-optimal neural hedging benchmarked against\\ Black--Scholes delta and no-trade-band strategies}
\author{Daniil Hutsol\\ \small Independent project $\cdot$ BSc student, SGH Warsaw School of Economics}
\date{October 2026}

\begin{document}
\maketitle

\begin{abstract}
We study the hedging of a short European call option under proportional transaction costs.
A recurrent neural-network policy, which observes the moneyness, the remaining time to
maturity and its current position, is trained end-to-end through a Monte Carlo market
simulator to minimise the Conditional Value-at-Risk (CVaR) of the terminal profit and loss.
Without costs, the network recovers the Black--Scholes delta from simulated data alone.
With 1\% proportional costs, it reduces $\CVaR_{95\%}$ out-of-sample by 41\% relative to
Black--Scholes delta hedging and by 8\% relative to the best tuned classical no-trade-band
strategy, a gap of roughly twelve times the seed-to-seed standard deviation. We document every
building block against theory and discuss the limitations of the simulated setting.
\end{abstract}

\tableofcontents
\newpage

%=====================================================================
\section{Introduction}
%=====================================================================

A dealer who sells an option is exposed to the movements of the underlying and reduces this
exposure by trading the underlying asset. In the frictionless Black--Scholes world the optimal
strategy is known in closed form: hold $\Delta = \partial C/\partial S$ shares and rebalance
continuously. In practice two assumptions fail. Rebalancing is discrete, which leaves a residual
hedging error, and every trade costs money, which makes frequent rebalancing expensive. Under
proportional transaction costs there is no exact closed-form optimal strategy; only asymptotic
results for small costs are available~\cite{whalley1997,leland1985}.

Deep hedging~\cite{buehler2019} replaces the analytical hedging rule by a neural network and
trains it directly on simulated market paths to minimise a chosen risk measure of the hedged
position. This report reproduces the core findings of that approach and extends the comparison
with classical benchmarks. Our contributions are:
\begin{enumerate}
  \item a fully validated pipeline: market simulator, pricing, hedging engine and training loop,
        each checked against analytical results;
  \item a CVaR-trained recurrent hedging network that accounts for its current position;
  \item a comparison not only with naive Black--Scholes delta hedging, but also with two
        classical cost-aware strategies (constant and Whalley--Wilmott no-trade bands) whose
        parameters are tuned on the training data;
  \item a robustness analysis over several training seeds.
\end{enumerate}

%=====================================================================
\section{Market model}
%=====================================================================

\subsection{Geometric Brownian motion}

The underlying price follows a geometric Brownian motion
\begin{equation}
  dS_t = \mu S_t\,dt + \sigma S_t\,dW_t ,
\end{equation}
where $\mu$ is the drift, $\sigma$ the volatility and $W_t$ a standard Brownian motion.
Its solution is
\begin{equation}
  S_t = S_0 \exp\!\Big( \big(\mu - \tfrac{1}{2}\sigma^2\big) t + \sigma W_t \Big),
\end{equation}
so $\log S_t$ is normally distributed with mean $\log S_0 + (\mu - \sigma^2/2)t$ and variance
$\sigma^2 t$. The correction $-\sigma^2/2$ ensures $\E[S_t] = S_0 e^{\mu t}$, since for
$X \sim \mathcal{N}(m, v)$ one has $\E[e^X] = e^{m + v/2}$.

\subsection{Exact discretisation}

On a grid $t_k = k\,\Delta t$, $\Delta t = T/N$, paths are generated exactly, without
discretisation error:
\begin{equation}
  S_{k+1} = S_k \exp\!\Big( \big(\mu - \tfrac{1}{2}\sigma^2\big)\Delta t + \sigma\sqrt{\Delta t}\, Z_{k+1} \Big),
  \qquad Z_{k+1} \overset{\text{iid}}{\sim} \mathcal{N}(0,1).
\end{equation}
The noise scales with $\sqrt{\Delta t}$ because variances of independent increments add up:
$N$ steps of variance $s^2$ must give the total variance $\sigma^2 T$, hence
$s^2 = \sigma^2 \Delta t$. Implementation-wise, the log-increments are summed with a cumulative
sum along the time axis, which vectorises the simulation over all paths.

\subsection{Parameters and validation}

Unless stated otherwise we use $S_0 = K = 100$, $\sigma = 0.2$, $\mu = r = 0$,
$T = 30/365$ and $N = 30$ daily rebalancing dates. Training uses $100{,}000$ paths
(seed 42); all reported results are computed on $100{,}000$ independent test paths (seed 1).

\begin{table}[H]
\centering
\begin{tabular}{lccc}
\toprule
Quantity & Theory & Simulation & Note \\
\midrule
$\E[S_T]$, $\mu = 0$        & 100.000 & 100.020 & difference $\approx 1$ standard error \\
$\E[S_T]$, $\mu = 0.5$      & 104.195 & 104.216 & checks the drift term \\
$\Var(\log S_T)$            & 0.003288 & 0.003258 & difference $\approx 2$ standard errors \\
\bottomrule
\end{tabular}
\caption{Validation of the simulator against closed-form moments.}
\end{table}

\maybefigure{../figures/geometric_brownian_simulation.png}{0.8}{Twenty simulated paths and the
5--95\% band of $100{,}000$ paths. The band widens as $\sqrt{t}$ and is slightly skewed upwards,
reflecting the log-normal distribution.}

%=====================================================================
\section{Pricing and the Black--Scholes delta}
%=====================================================================

With $r = 0$ the fair price of a European call is its expected payoff under the risk-neutral
measure, which coincides with the model with $\mu = 0$:
\begin{equation}
  C = \E\big[(S_T - K)^+\big] = S\,N(d_1) - K\,N(d_2),
\end{equation}
\begin{equation}
  d_1 = \frac{\ln(S/K) + \tfrac12 \sigma^2 \tau}{\sigma\sqrt{\tau}}, \qquad d_2 = d_1 - \sigma\sqrt{\tau},
\end{equation}
where $\tau = T - t$ is the remaining time to maturity and $N(\cdot)$ is the standard normal
distribution function. The Black--Scholes delta and gamma are
\begin{equation}
  \Delta = \frac{\partial C}{\partial S} = N(d_1), \qquad
  \Gamma = \frac{\partial^2 C}{\partial S^2} = \frac{N'(d_1)}{S\,\sigma\sqrt{\tau}} .
\end{equation}
Setting $r = 0$ is equivalent to working in discounted units and follows~\cite{buehler2019}.

For the at-the-money option of our setup $C = 2.2872$ and $\Delta = 0.5114$. The price matches the
Monte Carlo estimate $2.2896$ (difference below one quarter of a standard error), and the delta
matches the central finite difference $\big(C(S+h) - C(S-h)\big)/(2h)$ to eleven digits.

%=====================================================================
\section{The hedging problem}
%=====================================================================

\subsection{Profit and loss}

The dealer sells the call for the premium $p_0$ and holds $\delta_k$ units of the underlying on
$[t_k, t_{k+1})$, $k = 0, \dots, N-1$. With proportional costs $c$ per unit of traded notional,
the terminal profit and loss is
\begin{equation}\label{eq:pnl}
  \PnL = p_0 - (S_N - K)^+ + \sum_{k=0}^{N-1} \delta_k \,(S_{k+1} - S_k)
       - \sum_{k=0}^{N-1} c\, S_k \,\lvert \delta_k - \delta_{k-1} \rvert ,
  \qquad \delta_{-1} = 0 .
\end{equation}
The four terms are the premium, the option payoff, the gains of the hedge and the trading costs.
The turnover $\sum_k \lvert \delta_k - \delta_{k-1}\rvert$ measures the total number of shares traded.
We do not charge costs for unwinding the hedge at maturity, since the hedge ratio is then close to
0 or 1.

\subsection{Risk measures}

Let $L = -\PnL$ denote the loss. Besides the mean and the standard deviation we use
\begin{equation}
  \VaR_\alpha(L) = \inf\{\, \ell : \mathbb{P}(L \le \ell) \ge \alpha \,\}, \qquad
  \CVaR_\alpha(L) = \E\big[\,L \mid L \ge \VaR_\alpha(L)\,\big],
\end{equation}
the average loss in the worst $(1-\alpha)$ fraction of scenarios, with $\alpha = 0.95$.
Unlike the variance, CVaR only penalises the tail of losses and does not penalise gains.

\subsection{The Rockafellar--Uryasev representation}

Optimising CVaR directly requires sorting, which is inconvenient for gradient-based training.
Rockafellar and Uryasev~\cite{rockafellar2000} showed that
\begin{equation}\label{eq:ru}
  \CVaR_\alpha(L) = \min_{w \in \mathbb{R}} \Big\{\, w + \frac{1}{1-\alpha}\, \E\big[(L - w)^+\big] \Big\},
\end{equation}
with the minimum attained at $w = \VaR_\alpha(L)$. The objective inside the braces is an ordinary
expectation, which can be estimated on a mini-batch and differentiated. We treat $w$ as an
additional trainable parameter, so that minimising~\eqref{eq:ru} jointly over the network weights
and $w$ minimises the CVaR of the hedged position.

\subsection{Indifference premium}

CVaR is translation equivariant: increasing the premium by $\Delta p$ reduces CVaR by exactly
$\Delta p$. The smallest premium at which the hedged position has non-positive CVaR is therefore
\begin{equation}
  p^\ast = p_0 + \CVaR_\alpha(L).
\end{equation}
A better hedge thus translates directly into a lower price that can be quoted for the same tail risk.

%=====================================================================
\section{Classical benchmarks}
%=====================================================================

\subsection{Black--Scholes delta hedging}

The naive benchmark holds $\delta_k = \Delta(S_k, \tau_k)$ with $\tau_k = T - t_k$, regardless of
costs. Without costs, its discrete hedging error decays as $1/\sqrt{N}$: the error per step has
variance of order $\Delta t^2$ and the errors are approximately independent, so the total variance
is of order $N \Delta t^2 = T^2/N$. With costs, the expected change of the delta per step is of order
$\sqrt{\Delta t}$, so the turnover, and with it the trading cost, grows as
$N\sqrt{\Delta t} = \sqrt{NT} \propto \sqrt{N}$. Rebalancing more often therefore stops paying off.

Daily delta hedging already removes most of the risk of the short option: the P\&L standard
deviation falls from $3.46$ (unhedged) to $0.36$.

\maybefigure{../figures/pnl_hist_bs_vs_none.png}{0.8}{P\&L of the short call without hedging and with
daily BS delta hedging ($c = 0$, $N = 30$, log scale). Unhedged, the gain is capped at the premium
while losses are unbounded.}

\begin{table}[H]
\centering
\begin{tabular}{r cc cc cc}
\toprule
 & \multicolumn{2}{c}{$c = 0$} & \multicolumn{2}{c}{$c = 0.1\%$} & \multicolumn{2}{c}{$c = 1\%$} \\
\cmidrule(lr){2-3}\cmidrule(lr){4-5}\cmidrule(lr){6-7}
$N$ & mean & std & mean & std & mean & std \\
\midrule
5   & 0.000 & 0.839 & $-0.112$ & 0.849 & $-1.114$ & 0.953 \\
10  & 0.000 & 0.608 & $-0.147$ & 0.619 & $-1.451$ & 0.767 \\
30  & 0.000 & 0.358 & $-0.222$ & 0.373 & $-2.224$ & 0.756 \\
100 & 0.000 & 0.199 & $-0.367$ & 0.239 & $-3.679$ & 1.206 \\
\bottomrule
\end{tabular}
\caption{Black--Scholes delta hedging versus the number of rebalancing dates $N$. Without costs,
$\text{std}\cdot\sqrt{N} \approx 1.9$--$2.0$ is nearly constant. With 1\% costs, the mean loss at
$N=100$ exceeds the option premium ($2.29$) and the risk increases again.}
\end{table}

\maybefigure{../figures/hedging_error_vs_N.png}{0.6}{Discrete hedging error of BS delta hedging
without costs follows the $1/\sqrt{N}$ law.}

\maybefigure{../figures/bs_hedging_costs.png}{0.95}{Risk (left) and mean P\&L (right) of BS delta
hedging versus rebalancing frequency for three cost levels.}

\subsection{Constant no-trade band}

The classical response to proportional costs is a no-trade region around the delta. With
half-width $h$, the position is only adjusted when it leaves the band, and then only to the nearest
edge:
\begin{equation}
  \delta_k = \min\!\Big( \max\big( \delta_{k-1},\, \Delta_k - h \big),\, \Delta_k + h \Big).
\end{equation}
For $h = 0$ this is BS delta hedging. The width $h$ is tuned on the training paths to minimise
$\CVaR_{95\%}$; the optimum is $h = 0.11$.

\subsection{Whalley--Wilmott band}

Asymptotic analysis for small costs~\cite{whalley1997} suggests a state-dependent half-width that
scales with the option gamma:
\begin{equation}
  h_k = \kappa \,\big( c\, S_k\, \Gamma_k^2 \big)^{1/3},
\end{equation}
where all model constants (in particular risk aversion) are absorbed into the scale $\kappa$,
which is again tuned on the training paths ($\kappa = 0.7$).

%=====================================================================
\section{Deep hedging}
%=====================================================================

\subsection{Policy}

The hedging decision is a function of the observable state,
\begin{equation}
  \delta_k = f_\theta\big(x_{1,k},\, x_{2,k},\, \delta_{k-1}\big), \qquad
  x_{1,k} = \frac{\log(S_k/K)}{\sigma\sqrt{T}}, \qquad x_{2,k} = \frac{\tau_k}{T},
\end{equation}
where $f_\theta$ is a multilayer perceptron $3 \to 32 \to 32 \to 1$ with ReLU activations
(1{,}217 parameters). Inputs are normalised to order one; the first feature measures the moneyness
in units of the total standard deviation, as in $d_1$. The current position $\delta_{k-1}$ is
essential under costs: without it the policy cannot tell whether a trade is worth its cost.

Because the position at day $k$ is an input at day $k+1$, the policy is rolled out sequentially
over the $N$ dates, which makes the network recurrent. In the cost-free experiments the
position input is omitted ($2 \to 32 \to 32 \to 1$, 1{,}185 parameters) and all dates are
evaluated in a single forward pass.

\subsection{Training}

The policy is trained end-to-end through the simulator. At each iteration a mini-batch of
$10{,}000$ paths is drawn, the policy is rolled out, the P\&L~\eqref{eq:pnl} is computed with
differentiable tensor operations, and the loss~\eqref{eq:ru} (or the mean squared P\&L in the
first experiment) is backpropagated through all $N$ decisions to the network weights and to $w$.
We use Adam~\cite{kingma2015} with learning rate $10^{-3}$ for the network and $10^{-2}$ for $w$.

\subsection{Why not a tree-based model}

The problem has no labels: the optimal position at a given date is unknown, and the objective
is defined on the whole trajectory, where each decision feeds into the next state. Training
therefore requires a policy that is differentiable in its parameters, so that the gradient of
the risk measure can be propagated through the simulator. Tree ensembles are piecewise constant
and not differentiable in their parameters, so they do not fit this end-to-end formulation.
The approach is closest to direct policy optimisation in reinforcement learning, with the
difference that here the environment itself is differentiable and gradients are exact.

%=====================================================================
\section{Experiments}
%=====================================================================

\subsection{Without costs: recovering Black--Scholes}

Trained on the mean squared P\&L, the network learns positions that differ from the Black--Scholes
delta by $0.0086$ on average, out-of-sample. Its P\&L standard deviation is $0.369$ versus $0.358$
for BS delta, and the results on training and test paths coincide, so the network has learned
the hedging rule rather than memorising paths.

\maybefigure{../figures/net_vs_bs_delta.png}{0.7}{Learned hedge ratio versus the Black--Scholes
delta for several dates, without transaction costs.}

\subsection{Without costs: CVaR versus variance}

\begin{table}[H]
\centering
\begin{tabular}{lccc}
\toprule
Strategy & std & $\CVaR_{95\%}$ \\
\midrule
BS delta            & \textbf{0.358} & 0.833 \\
Network, MSE loss   & 0.369 & 0.844 \\
Network, CVaR loss  & 0.409 & \textbf{0.739} \\
\bottomrule
\end{tabular}
\caption{Out-of-sample comparison without transaction costs.}
\end{table}

Each strategy is best on the criterion it targets. The CVaR-trained network shortens the left
tail by 11\% at the cost of a wider distribution: the Black--Scholes delta is close to
variance-optimal but not tail-optimal.

\maybefigure{../figures/pnl_bs_vs_cvar_net.png}{0.8}{P\&L distributions (log scale) of BS delta and
the CVaR-trained network; dashed lines mark $-\VaR_{95\%}$.}

\subsection{With 1\% transaction costs}

\begin{table}[H]
\centering
\begin{tabular}{lcccc}
\toprule
Strategy & mean P\&L & std & $\CVaR_{95\%}$ & turnover \\
\midrule
No hedge                               & $-0.02$ & 3.47 & 10.06 & 0.00 \\
BS delta                               & $-2.22$ & 0.75 & 4.06 & 2.22 \\
Constant band, $h = 0.11$              & $-1.04$ & 0.71 & 2.61 & 1.03 \\
Whalley--Wilmott band, $\kappa = 0.7$  & $-1.04$ & 0.78 & 2.76 & 1.03 \\
Network (5 seeds)                      & $-1.22 \pm 0.02$ & $\mathbf{0.65 \pm 0.01}$ & $\mathbf{2.40 \pm 0.02}$ & $\approx 1.3$ \\
\bottomrule
\end{tabular}
\caption{Out-of-sample results with $c = 1\%$. Band parameters are tuned on the training paths;
network results are mean $\pm$ standard deviation over five training seeds.}
\end{table}

\maybefigure{../figures/benchmark_comparison.png}{0.95}{Left: tuning of the constant band width on
the training paths. Right: out-of-sample $\CVaR_{95\%}$ of all strategies; the error bar shows the
seed-to-seed standard deviation of the network.}

The main observations are:
\begin{itemize}
  \item Most of the improvement over BS delta comes from the no-trade principle itself: a constant
        band with a single tuned parameter reduces $\CVaR_{95\%}$ by 36\%.
  \item The network improves on the best classical benchmark by a further 8\% in $\CVaR_{95\%}$ and
        9\% in standard deviation. The gap of $0.21$ is about twelve times the seed-to-seed standard
        deviation, and every individual seed beats the band.
  \item The network accepts a slightly larger mean loss than the band: it trades somewhat more
        to cut the tail, which is consistent with its objective.
  \item The Whalley--Wilmott band underperforms the constant band. It is derived for small costs
        and utility maximisation; under a CVaR objective with 1\% costs its widening near expiry,
        where the gamma of an at-the-money option explodes, leaves the position under-hedged exactly
        when the tail outcome is decided.
  \item The indifference premium for the same $\CVaR_{95\%}$ is $2.29 + 2.40 \approx 4.70$ for the
        network versus $2.29 + 4.06 \approx 6.35$ for BS delta, so the option could be quoted about
        26\% cheaper.
\end{itemize}

\subsection{What the network learned}

Fixing the date and the price and varying only the current position reveals the learned trading
rule. Unlike the BS delta, which jumps to the same target from any starting position, the learned
rule has a slope well below one in the current position: the network moves only part of the way
towards its target and spreads the adjustment over several days, which saves transaction costs.
The fixed point of the rule, where no trade occurs, lies close to the BS delta for each price level.
This partial rebalancing is a smooth approximation of the band structure of the theoretical
solution rather than an exact flat no-trade region. Regions of the state space that are rarely
visited during training, such as a zero position at $S = 105$ in mid-life, should be interpreted
with caution.

\maybefigure{../figures/no_trade_band.png}{0.6}{Learned trading rule at day 15 with $c = 1\%$:
new position versus current position for three price levels. The diagonal corresponds to no
trade; dotted lines show the BS delta.}

%=====================================================================
\section{Limitations and model risk}
%=====================================================================

\begin{itemize}
  \item \textbf{Market dynamics.} Training and testing use geometric Brownian motion, which has
        neither fat tails nor volatility clustering. Results on real markets may differ.
  \item \textbf{Premium.} The option is sold at the frictionless Black--Scholes price rather than
        at a market quote. Relative comparisons between strategies are unaffected, since the premium
        is identical for all of them, but absolute profitability is not assessed.
  \item \textbf{Frictions.} Costs are proportional and constant; there is no market impact and no
        cost for unwinding at maturity.
  \item \textbf{Scope.} A single at-the-money option with one maturity, hedged in the underlying only.
\end{itemize}

%=====================================================================
\section{Next steps}
%=====================================================================

\begin{itemize}
  \item Replace GBM by a block bootstrap of historical S\&P~500 returns (2005--2017), which preserves
        fat tails and volatility clustering, and test on realised paths from 2018--2026.
  \item Add an ex-ante volatility feature to the policy (trailing realised volatility, a GARCH(1,1)
        forecast with parameters estimated on the training period only, or the VIX), and report
        results by volatility regime and in historical stress windows.
  \item Forecast VaR and ES of the hedged book with a neural quantile model and backtest it
        (exceedance rate, Kupiec test) against historical simulation, linear quantile regression and
        gradient boosting.
\end{itemize}

%=====================================================================
\begin{thebibliography}{9}

\bibitem{black1973}
F.~Black and M.~Scholes.
\newblock The pricing of options and corporate liabilities.
\newblock \emph{Journal of Political Economy}, 81(3):637--654, 1973.

\bibitem{buehler2019}
H.~Buehler, L.~Gonon, J.~Teichmann, and B.~Wood.
\newblock Deep hedging.
\newblock \emph{Quantitative Finance}, 19(8):1271--1291, 2019.

\bibitem{kingma2015}
D.~P.~Kingma and J.~Ba.
\newblock Adam: A method for stochastic optimization.
\newblock In \emph{International Conference on Learning Representations}, 2015.

\bibitem{leland1985}
H.~E.~Leland.
\newblock Option pricing and replication with transactions costs.
\newblock \emph{The Journal of Finance}, 40(5):1283--1301, 1985.

\bibitem{rockafellar2000}
R.~T.~Rockafellar and S.~Uryasev.
\newblock Optimization of conditional value-at-risk.
\newblock \emph{Journal of Risk}, 2(3):21--41, 2000.

\bibitem{whalley1997}
A.~E.~Whalley and P.~Wilmott.
\newblock An asymptotic analysis of an optimal hedging model for option pricing with transaction costs.
\newblock \emph{Mathematical Finance}, 7(3):307--324, 1997.

\end{thebibliography}

\end{document}
